%%%%%%%%%%%%%%%%%%%%%%%%%%%%%%%%%%%%%%%%%%%%%%%%%%%%%%%%%%%%%%%%%%%%%%%%%%%%%%
%
% 06_coherence.tex
%
%%%%%%%%%%%%%%%%%%%%%%%%%%%%%%%%%%%%%%%%%%%%%%%%%%%%%%%%%%%%%%%%%%%%%%%%%%%%%%

\begin{mosbox}{5. Sheaf Coherence and Global Consistency}

\BodyFont

Once local information has been attached to the simplicial complex through
a cellular sheaf, the next question is whether these local assignments are
globally compatible.

MOS measures this compatibility through the sheaf Laplacian, which
quantifies disagreement between neighboring local sections.

Rather than asking whether individual concepts are correct, the objective
is to determine whether the entire assignment is internally consistent.

\end{mosbox}

%%%%%%%%%%%%%%%%%%%%%%%%%%%%%%%%%%%%%%%%%%%%%%%%%%%%%%%%%%%%%%%%%%%%%%%%%%%%%%
%% SHEAF LAPLACIAN
%%%%%%%%%%%%%%%%%%%%%%%%%%%%%%%%%%%%%%%%%%%%%%%%%%%%%%%%%%%%%%%%%%%%%%%%%%%%%%

\begin{eqbox}

\[
\boxed{
L=\delta^{T}\delta
}
\]

\vspace{3mm}

\BodyFont

Here

\[
\delta
\]

denotes the sheaf coboundary operator.

The associated Laplacian measures incompatibility introduced by the local
restriction maps.

\end{eqbox}

%%%%%%%%%%%%%%%%%%%%%%%%%%%%%%%%%%%%%%%%%%%%%%%%%%%%%%%%%%%%%%%%%%%%%%%%%%%%%%
%% ENERGY
%%%%%%%%%%%%%%%%%%%%%%%%%%%%%%%%%%%%%%%%%%%%%%%%%%%%%%%%%%%%%%%%%%%%%%%%%%%%%%

\vspace{5mm}

\begin{mosbox}{Coherence Functional}

\BodyFont

Given an assignment

\[
x\in C^{0}(C,\mathcal F),
\]

its disagreement energy is defined by

\[
\boxed{
\omega(x)
=
x^{T}Lx
=
\|\delta x\|^{2}.
}
\]

This functional has a natural interpretation:

\begin{center}

\begin{tabular}{p{5cm}p{10cm}}

$\omega(x)=0$

&

Perfect global compatibility.

\\[4mm]

Small $\omega(x)$

&

Minor local disagreement.

\\[4mm]

Large $\omega(x)$

&

Substantial inconsistency between neighboring assignments.

\end{tabular}

\end{center}

The quantity therefore serves as a global measure of structural coherence.

\end{mosbox}

%%%%%%%%%%%%%%%%%%%%%%%%%%%%%%%%%%%%%%%%%%%%%%%%%%%%%%%%%%%%%%%%%%%%%%%%%%%%%%
%% GEOMETRIC INTERPRETATION
%%%%%%%%%%%%%%%%%%%%%%%%%%%%%%%%%%%%%%%%%%%%%%%%%%%%%%%%%%%%%%%%%%%%%%%%%%%%%%

\begin{center}

\begin{tikzpicture}[
  scale=1.2,
  healthy/.style={circle, fill=MOSGold, draw=white, line width=1pt, minimum size=10pt, inner sep=0pt},
  discord/.style={circle, fill=MOSRed, draw=white, line width=1pt, minimum size=10pt, inner sep=0pt},
  hedge/.style={draw=MOSGold, line width=2pt},
  dedge/.style={draw=MOSRed, line width=2pt},
  broken/.style={draw=MOSRed, line width=2pt, decorate, decoration={zigzag, segment length=4mm, amplitude=1.5mm}}
]

% Healthy side
\node[white, font=\Large\bfseries] at (-4, 3) {healthy};

\node[healthy] (h1) at (-5.5, 1.5) {};
\node[healthy] (h2) at (-4, 1.5) {};
\node[healthy] (h3) at (-2.5, 1.5) {};
\node[healthy] (h4) at (-5.5, 0) {};
\node[healthy] (h5) at (-4, 0) {};

\draw[hedge] (h1) -- (h2) -- (h3);
\draw[hedge] (h1) -- (h4);
\draw[hedge] (h2) -- (h5);
\draw[hedge] (h4) -- (h5);

% Rho
\node[white, font=\Huge\bfseries] at (0, 1.5) {$\rho$};

% Discord side
\node[white, font=\Large\bfseries] at (4, 3) {discord};

\node[discord] (d1) at (2.5, 1.5) {};
\node[discord] (d2) at (4, 1.5) {};
\node[discord] (d3) at (5.5, 1.5) {};
\node[discord] (d4) at (2.5, 0) {};

\draw[dedge] (d1) -- (d2);
\draw[broken] (d2) -- (d3) node[midway, fill=MOSBlack, font=\Huge, text=white, inner sep=1pt] {$\times$};
\draw[dedge] (d1) -- (d4);

\end{tikzpicture}

\end{center}

%%%%%%%%%%%%%%%%%%%%%%%%%%%%%%%%%%%%%%%%%%%%%%%%%%%%%%%%%%%%%%%%%%%%%%%%%%%%%%
%% KERNEL
%%%%%%%%%%%%%%%%%%%%%%%%%%%%%%%%%%%%%%%%%%%%%%%%%%%%%%%%%%%%%%%%%%%%%%%%%%%%%%

\begin{highlightbox}

\SectionTitle{Global Sections}

\BodyFont

A fundamental result used throughout MOS is

\[
\boxed{
H^{0}(C,\mathcal F)
=
\ker(L).
}
\]

Consequently,

global sections are precisely those assignments whose disagreement energy
vanishes.

Equivalently,

\[
\omega(x)=0
\Longleftrightarrow
x\in H^{0}(C,\mathcal F).
\]

This provides a direct computational characterization of coherent memory
states.

\end{highlightbox}

%%%%%%%%%%%%%%%%%%%%%%%%%%%%%%%%%%%%%%%%%%%%%%%%%%%%%%%%%%%%%%%%%%%%%%%%%%%%%%
%% NORMALIZED DISCORD
%%%%%%%%%%%%%%%%%%%%%%%%%%%%%%%%%%%%%%%%%%%%%%%%%%%%%%%%%%%%%%%%%%%%%%%%%%%%%%

\begin{mosbox}{Normalized Discord}

\BodyFont

For comparisons across different systems and scales, MOS introduces the
dimensionless normalized discord

\[
\boxed{
\rho
=
\frac{\omega(x)}
{B\|x\|^{2}},
\qquad
0\le\rho\le1,
}
\]

where

\[
B
\]

is an appropriate normalization constant.

The normalized quantity removes dependence on the overall scale of the
assignment, allowing coherence values to be compared across distinct
memory structures.

\end{mosbox}

%%%%%%%%%%%%%%%%%%%%%%%%%%%%%%%%%%%%%%%%%%%%%%%%%%%%%%%%%%%%%%%%%%%%%%%%%%%%%%
%% SCALE
%%%%%%%%%%%%%%%%%%%%%%%%%%%%%%%%%%%%%%%%%%%%%%%%%%%%%%%%%%%%%%%%%%%%%%%%%%%%%%

\begin{center}

\begin{tikzpicture}

\draw[
MOSGray,
line width=3pt
]
(0,0)
--
(12,0);

\foreach \x/\lab in
{
0/{0},
3/{0.25},
6/{0.5},
9/{0.75},
12/{1}
}
{
\draw[
MOSGray,
line width=2pt
]
(\x,-.15)
--
(\x,.15);

\node[below]
at
(\x,-.15)
{$\lab$};
}

\fill[MOSGold]
(0,0)
circle(5pt);

\fill[MOSGold]
(12,0)
circle(5pt);

\node[above]
at
(0,.3)
{\small Perfect Coherence};

\node[above]
at
(12,.3)
{\small Maximum Discord};

\end{tikzpicture}

\end{center}

%%%%%%%%%%%%%%%%%%%%%%%%%%%%%%%%%%%%%%%%%%%%%%%%%%%%%%%%%%%%%%%%%%%%%%%%%%%%%%
%% REMARK
%%%%%%%%%%%%%%%%%%%%%%%%%%%%%%%%%%%%%%%%%%%%%%%%%%%%%%%%%%%%%%%%%%%%%%%%%%%%%%

\begin{remarkbox}

\BodyFont

The coherence functional is one of the central diagnostic tools of MOS.

It transforms the abstract notion of compatibility into a computable scalar
quantity while preserving its geometric interpretation through the sheaf
Laplacian.

Subsequent sections build upon this measure to develop semantic geometry,
diagnostics and adaptive structural dynamics.

\end{remarkbox}